\section{Characterizing finite and finitely presented modules}
\label{section-colim-and-hom}

\noindent
Given a module $N$ over a ring $R$, you can characterize whether or not
$N$ is a finite module or a finitely presented module
in terms of the functor $\Hom_R(N, -)$.

\begin{lemma}
\label{lemma-characterize-finite-module-hom}
Let $R$ be a ring. Let $N$ be an $R$-module. The following are equivalent
\begin{enumerate}
\item $N$ is a finite $R$-module,
\item for any filtered colimit $M = \colim M_i$ of $R$-modules the map
$\colim \Hom_R(N, M_i) \to \Hom_R(N, M)$ is injective.
\end{enumerate}
\end{lemma}

\begin{proof}
Assume (1) and choose generators $x_1, \ldots, x_m$ for $N$.
If $N \to M_i$ is a module map and the composition
$N \to M_i \to M$ is zero, then because $M = \colim_{i' \geq i} M_{i'}$
for each $j \in \{1, \ldots, m\}$ we can find a $i' \geq i$ such that
$x_j$ maps to zero in $M_{i'}$. Since there are finitely many
$x_j$ we can find a single $i'$ which works for all of them.
Then the composition $N \to M_i \to M_{i'}$ is zero and we conclude
the map is injective, i.e., part (2) holds.

\medskip\noindent
Assume (2). For a finite subset $E \subset N$ denote $N_E \subset N$
the $R$-submodule generated by the elements of $E$. Then
$0 = \colim N/N_E$ is a filtered colimit. Hence we see that
$\text{id} : N \to N$ maps into $N_E$ for some $E$, i.e., $N$
is finitely generated.
\end{proof}

\noindent
For purposes of reference, we define what it means to have a relation
between elements of a module.

\begin{definition}
\label{definition-relation}
Let $R$ be a ring. Let $M$ be an $R$-module.
Let $n \geq 0$ and $x_i \in M$ for $i = 1, \ldots, n$.
A {\it relation} between $x_1, \ldots, x_n$ in $M$ is a
sequence of elements $f_1, \ldots, f_n \in R$ such that
$\sum_{i = 1, \ldots, n} f_i x_i = 0$.
\end{definition}

\begin{lemma}
\label{lemma-module-colimit-fp}
Let $R$ be a ring and let $M$ be an $R$-module.
Then $M$ is the colimit of a directed system
$(M_i, \mu_{ij})$ of $R$-modules
with all $M_i$ finitely presented $R$-modules.
\end{lemma}

\begin{proof}
Consider any finite subset $S \subset M$ and any finite
collection of relations $E$ among the elements
of $S$. So each $s \in S$ corresponds to $x_s \in M$ and
each $e \in E$ consists of a vector
of elements $f_{e, s} \in R$ such that $\sum f_{e, s} x_s = 0$.
Let $M_{S, E}$ be the cokernel of the map
$$
R^{\# E} \longrightarrow R^{\# S}, \quad
(g_e)_{e\in E} \longmapsto (\sum g_e f_{e, s})_{s\in S}.
$$
There are canonical maps $M_{S, E} \to M$.
If $S \subset S'$ and if the elements of
$E$ correspond, via this map, to relations
in $E'$, then there is an obvious map
$M_{S, E} \to M_{S', E'}$ commuting with the
maps to $M$. Let $I$ be the set of pairs
$(S, E)$ with ordering by inclusion as above.
It is clear that the colimit of this directed system is $M$.
\end{proof}

\begin{lemma}
\label{lemma-characterize-finitely-presented-module-hom}
Let $R$ be a ring. Let $N$ be an $R$-module. The following are equivalent
\begin{enumerate}
\item $N$ is a finitely presented $R$-module,
\item for any filtered colimit $M = \colim M_i$ of $R$-modules the map
$\colim \Hom_R(N, M_i) \to \Hom_R(N, M)$ is bijective.
\end{enumerate}
\end{lemma}

\begin{proof}
Assume (1) and choose an exact sequence $F_{-1} \to F_0 \to N \to 0$
with $F_i$ finite free. Then we have an exact sequence
$$
0 \to \Hom_R(N, M) \to \Hom_R(F_0, M) \to \Hom_R(F_{-1}, M)
$$
functorial in the $R$-module $M$. The functors $\Hom_R(F_i, M)$ commute
with filtered colimits as $\Hom_R(R^{\oplus n}, M) = M^{\oplus n}$.
Since filtered colimits are exact
(Lemma \ref{lemma-directed-colimit-exact})
we see that (2) holds.

\medskip\noindent
Assume (2). By Lemma \ref{lemma-module-colimit-fp}
we can write $N = \colim N_i$ as a filtered
colimit such that $N_i$ is of finite presentation for all $i$.
Thus $\text{id}_N$ factors through $N_i$ for some $i$.
This means that $N$ is a direct summand of a finitely
presented $R$-module (namely $N_i$) and hence finitely presented.
\end{proof}














